\documentclass[12pt, a4paper]{article}

\usepackage[utf8]{inputenc}
\usepackage[T1]{fontenc}
\usepackage[brazil]{babel}
\usepackage{geometry}
\usepackage{listings}
\usepackage{xcolor}
\usepackage{hyperref}
\usepackage{graphicx}
\usepackage{titlesec}
\usepackage{fancyhdr}
\usepackage{setspace}
\usepackage{indentfirst}
\usepackage{booktabs}

\geometry{top=3cm, bottom=2cm, left=3cm, right=2cm}

% Configuração de cores e estilo para código
\definecolor{codebg}{rgb}{0.95,0.95,0.95}
\definecolor{codegreen}{rgb}{0,0.5,0}
\definecolor{codegray}{rgb}{0.5,0.5,0.5}
\definecolor{codepurple}{rgb}{0.4,0,0.6}
\definecolor{codeblue}{rgb}{0,0,0.8}

\lstdefinestyle{python}{
    backgroundcolor=\color{codebg},
    commentstyle=\color{codegreen},
    keywordstyle=\color{codeblue}\bfseries,
    numberstyle=\tiny\color{codegray},
    stringstyle=\color{codepurple},
    basicstyle=\ttfamily\footnotesize,
    breaklines=true,
    captionpos=b,
    numbers=left,
    numbersep=5pt,
    showspaces=false,
    frame=single,
    rulecolor=\color{black!30},
    tabsize=4
}

\lstdefinestyle{minilang}{
    backgroundcolor=\color{codebg},
    keywordstyle=\color{codeblue}\bfseries,
    basicstyle=\ttfamily\footnotesize,
    breaklines=true,
    captionpos=b,
    numbers=left,
    numbersep=5pt,
    frame=single,
    rulecolor=\color{black!30},
    morekeywords={LET, PRINT, IF, ELIF, ELSE, FOR, WHILE, END, TO},
    tabsize=4
}

\pagestyle{fancy}
\fancyhf{}
\rhead{Compiladores -- AV2}
\lhead{Mini Compiler -- Minilang}
\cfoot{\thepage}

\onehalfspacing

\begin{document}

% ==========================
% CAPA
% ==========================
\begin{titlepage}
    \centering
    \vspace*{2cm}
    {\large \textbf{UNINASSAU}}\\[0.5cm]
    {\large Disciplina de Compiladores}\\[3cm]

    {\Huge \textbf{Mini Compiler}}\\[0.5cm]
    {\Large Minilang -- Compilador Educacional em Python}\\[3cm]

    {\large \textbf{Integrantes do Grupo:}}\\[0.3cm]
    {\large Ryan Nunes, Bryan Samuel, Anderson Djalma e Gabriel Barros}\\[1cm]

    {\large \textbf{Professor:} Anderlan}\\[1cm]
    {\large Maceió - AL}\\[0.3cm]
    {\large \today}
\end{titlepage}

% ==========================
% SUMÁRIO
% ==========================
\tableofcontents
\newpage

% ==========================
% 1. INTRODUÇÃO
% ==========================
\section{Introdução}

Este relatório documenta o desenvolvimento do \textbf{Mini Compiler}, um compilador educacional
criado em Python para a disciplina de Compiladores. O projeto tem como objetivo implementar
as etapas fundamentais de um compilador real: análise léxica, análise sintática, geração de
Árvore Sintática Abstrata (AST) e geração de código.

A linguagem desenvolvida, chamada \textbf{Minilang}, é uma linguagem simples com sintaxe
em português/inglês que compila para código Python. Ela suporta variáveis, operações
matemáticas, estruturas condicionais e de repetição.

% ==========================
% 2. ESTRUTURA DO COMPILADOR
% ==========================
\section{Estrutura do Compilador}

Um compilador é um programa que traduz código escrito em uma linguagem (linguagem fonte)
para outra linguagem (linguagem destino). O Mini Compiler é composto por quatro módulos
principais, conforme ilustrado abaixo:

\begin{center}
\begin{tabular}{ccccc}
\texttt{programa.mini} & $\rightarrow$ & \textbf{Lexer} & $\rightarrow$ & \textbf{Tokens} \\
& & & & $\downarrow$ \\
\texttt{output.py} & $\leftarrow$ & \textbf{Generator} & $\leftarrow$ & \textbf{Parser (AST)} \\
\end{tabular}
\end{center}

\subsection{Lexer (Análise Léxica)}

O \textbf{Lexer} é a primeira etapa do compilador. Ele lê o código fonte caractere por caractere
e agrupa em unidades chamadas \textbf{tokens}. Um token é o menor elemento com significado
em uma linguagem, como uma palavra reservada, um número ou um operador.

O Lexer do Minilang foi implementado usando a biblioteca \textbf{PLY (Python Lex-Yacc)}.
Os tokens reconhecidos são:

\begin{table}[h]
\centering
\begin{tabular}{@{}ll@{}}
\toprule
\textbf{Token} & \textbf{Descrição} \\
\midrule
\texttt{LET} & Declaração de variável \\
\texttt{PRINT} & Impressão de valor \\
\texttt{IF / ELIF / ELSE} & Estrutura condicional \\
\texttt{FOR} & Laço de repetição com contador \\
\texttt{WHILE} & Laço de repetição com condição \\
\texttt{END} & Fechamento de bloco \\
\texttt{TO} & Intervalo do FOR \\
\texttt{NUMBER} & Número inteiro ou decimal \\
\texttt{STRING} & Texto entre aspas \\
\texttt{ID} & Identificador (nome de variável) \\
\texttt{EQUALS / EQUAL\_EQUAL} & Atribuição (=) e comparação (==) \\
\texttt{GREATER / LESS} & Operadores \textgreater{} e \textless{} \\
\texttt{GREATER\_EQUAL / LESS\_EQUAL} & Operadores \textgreater{}= e \textless{}= \\
\texttt{NOT\_EQUAL} & Operador != \\
\texttt{PLUS / MINUS / TIMES / DIVIDE} & Operadores aritméticos \\
\bottomrule
\end{tabular}
\caption{Tokens da linguagem Minilang}
\end{table}

Exemplo de como o Lexer tokeniza o código:

\begin{lstlisting}[style=minilang, caption={Entrada no Lexer}]
LET x = 10
IF x > 5
PRINT "Maior"
END
\end{lstlisting}

\begin{lstlisting}[style=python, caption={Tokens gerados pelo Lexer}]
LexToken(LET, 'LET', 1, 0)
LexToken(ID, 'x', 1, 4)
LexToken(EQUALS, '=', 1, 6)
LexToken(NUMBER, 10, 1, 8)
LexToken(IF, 'IF', 2, 11)
LexToken(ID, 'x', 2, 14)
LexToken(GREATER, '>', 2, 16)
LexToken(NUMBER, 5, 2, 18)
LexToken(PRINT, 'PRINT', 3, 20)
LexToken(STRING, 'Maior', 3, 26)
LexToken(END, 'END', 4, 33)
\end{lstlisting}

\subsection{Parser (Análise Sintática)}

O \textbf{Parser} recebe a lista de tokens do Lexer e verifica se eles formam uma estrutura
gramaticalmente válida, de acordo com as regras da linguagem. Ele utiliza o método
\textbf{LALR (Look-Ahead LR)}, implementado pelo PLY, que é o mesmo método utilizado
em compiladores profissionais.

O resultado do Parser é uma \textbf{Árvore Sintática Abstrata (AST)}, que representa
a estrutura hierárquica do programa.

Exemplo de AST gerada para o código acima:

\begin{lstlisting}[style=python, caption={AST gerada pelo Parser}]
[
  {'type': 'LET', 'name': 'x', 'value': 10},
  {'type': 'IF',
   'left': 'x', 'operator': '>', 'right': 5,
   'body': [{'type': 'PRINT', 'value': 'Maior',
             'value_type': 'STRING'}],
   'elifs': [], 'else': None}
]
\end{lstlisting}

\subsection{Generator (Geração de Código)}

O \textbf{Generator} percorre a AST e traduz cada nó para código Python equivalente.
Ele utiliza recursão para tratar comandos aninhados (por exemplo, um \texttt{LET}
dentro de um \texttt{IF}), garantindo que a indentação do Python seja gerada corretamente.

\begin{lstlisting}[style=python, caption={Código Python gerado}]
x = 10
if x > 5:
    print("Maior")
\end{lstlisting}

% ==========================
% 3. SINTAXE DA LINGUAGEM
% ==========================
\section{Sintaxe da Linguagem Minilang}

\subsection{Variáveis e Strings}
\begin{lstlisting}[style=minilang]
LET nome = "Brian"
LET idade = 22
LET soma = idade + 10
\end{lstlisting}

\subsection{Operações Matemáticas}
\begin{lstlisting}[style=minilang]
LET a = 10
LET b = 3
LET resultado = a * b
PRINT resultado
\end{lstlisting}

\subsection{Estrutura Condicional}
\begin{lstlisting}[style=minilang]
LET nota = 7

IF nota > 7
PRINT "Aprovado"

ELIF nota == 7
PRINT "Passou na media"

ELSE
PRINT "Reprovado"

END
\end{lstlisting}

\subsection{Laço FOR}
\begin{lstlisting}[style=minilang]
FOR i = 1 TO 5
PRINT i
END
\end{lstlisting}

\subsection{Laço WHILE}
\begin{lstlisting}[style=minilang]
LET x = 1
WHILE x < 6
PRINT x
LET x = x + 1
END
\end{lstlisting}

% ==========================
% 4. TECNOLOGIAS UTILIZADAS
% ==========================
\section{Tecnologias Utilizadas}

\begin{itemize}
    \item \textbf{Python 3.13} -- Linguagem de implementação do compilador
    \item \textbf{PLY (Python Lex-Yacc)} -- Biblioteca para construção do Lexer e Parser LALR
    \item \textbf{Git e GitHub} -- Controle de versão e colaboração
    \item \textbf{VS Code} -- Ambiente de desenvolvimento
    \item \textbf{MikTeX + TeXMaker} -- Compilação e edição deste relatório em \LaTeX{}
\end{itemize}

% ==========================
% 5. USO DE IA
% ==========================
\section{Uso de Inteligência Artificial}

Conforme permitido pelas regras da AV2, o grupo utilizou Inteligência Artificial
(Claude -- Anthropic) como apoio no desenvolvimento do projeto. A seguir estão
os principais usos e os erros encontrados no processo.

\subsection{Como a IA foi utilizada}

\begin{itemize}
    \item Revisão e melhoria do \texttt{generator.py}
    \item Correção de erros de importação no Python
    \item Criação do arquivo \texttt{programa.mini} com exemplos completos
    \item Colaboração deste relatório em \LaTeX{}
\end{itemize}

\subsection{Erros encontrados durante o desenvolvimento}

\begin{enumerate}
    \item \textbf{ImportError: cannot import name 'lexer' from 'lexer'} \\
          O arquivo \texttt{lexer.py} estava incompleto -- faltava a linha \texttt{lexer = lex.lex()}
          no final, que cria e exporta o objeto do lexer.

    \item \textbf{ImportError: cannot import name 'parser' from 'parser'} \\
          O Python possui um módulo nativo chamado \texttt{parser}, que conflitava com o arquivo
          \texttt{parser.py} do projeto. A solução foi renomear o arquivo para
          \texttt{minilang\_parser.py} e atualizar os imports.

\end{enumerate}

% ==========================
% 6. COMO EXECUTAR
% ==========================
\section{Como Executar o Projeto}

\subsection{Pré-requisitos}
\begin{lstlisting}[language=bash]
pip install ply
\end{lstlisting}

\subsection{Execução}
\begin{enumerate}
    \item Escreva o código Minilang no arquivo \texttt{programa.mini}
    \item Execute o compilador:
\begin{lstlisting}[language=bash]
python main.py
\end{lstlisting}
    \item O código Python gerado será salvo em \texttt{output.py}
    \item Para executar o resultado:
\begin{lstlisting}[language=bash]
python output.py
\end{lstlisting}
\end{enumerate}

% ==========================
% 7. CONCLUSÃO
% ==========================
\section{Conclusão}

O desenvolvimento do Mini Compiler permitiu ao grupo compreender na prática as etapas
fundamentais de um compilador: análise léxica, análise sintática, construção da AST e
geração de código. A utilização da biblioteca PLY facilitou a implementação do Lexer e
do Parser LALR, enquanto a refatoração do Generator com recursão garantiu suporte a
estruturas aninhadas.

O projeto atingiu todos os requisitos da AV2: possui condicional (IF/ELIF/ELSE),
laços de repetição (FOR e WHILE), versionamento no GitHub, e este relatório
elaborado no editor tipográfico \LaTeX{}.

\end{document}
